\chapter{Algebraic QFT: Local Nets and States}
\label{ch:volume-iv-aqft-local-nets-states}
\ConventionsLorentzian


\section{Local Observable Nets}

Let \(M=\mathbb M^D\) be Minkowski spacetime, and let
\(\mathcal K(M)\) be a chosen class of open spacetime regions, for example
double cones or bounded causally complete regions.  The connected Poincare
group is denoted
\[
  \Poinc=\mathbb R^D\rtimes SO^+(1,D-1).
\]
A local net of observables assigns to each
\[
  \mathcal O\in\mathcal K(M)
\]
a unital \(C^*\)-algebra \(\Obs(\mathcal O)\).  Throughout this chapter
\(\Obs(\mathcal O)\) denotes the abstract \(C^*\)-local algebra generated by
observables localized inside \(\mathcal O\), before a Hilbert-space
representation has been chosen.  Von Neumann local algebras are introduced only
after representing the \(C^*\)-net in a state.

Covariance is specified by a group homomorphism
\[
  g\longmapsto \alpha_g
\]
from \(\Poinc\) to automorphisms of the net, equivalently to
automorphisms of the quasilocal algebra after the inductive limit below has
been formed.

The basic net properties are:
\[
\begin{array}{ll}
  \text{isotony:} &
  \mathcal O_1\subset \mathcal O_2
  \Rightarrow \Obs(\mathcal O_1)\subset\Obs(\mathcal O_2),\\[0.25em]
  \text{Einstein causality:} &
  \mathcal O_1\perp \mathcal O_2\Rightarrow [A,B]=0
  \text{ for }A\in\Obs(\mathcal O_1),\ B\in\Obs(\mathcal O_2),\\[0.25em]
  \text{covariance:} &
  \alpha_g(\Obs(\mathcal O))=\Obs(g\mathcal O),
  \qquad g\in \Poinc,\\[0.25em]
  \text{time-slice:} &
  \Obs(N)=\Obs(D(N))\text{ when }N\text{ contains a Cauchy surface for }D(N).
\end{array}
\]
Here \(\mathcal O_1\perp\mathcal O_2\) means spacelike separation,
and \(D(N)\) is the domain of dependence of \(N\).  Additivity is often
imposed:
\[
  \mathcal O=\bigcup_i \mathcal O_i
  \quad\Longrightarrow\quad
  \Obs(\mathcal O)\text{ is generated by the }\Obs(\mathcal O_i).
\]

The quasilocal algebra is the \(C^*\)-inductive limit
\[
  \Obs_{\rm qloc}
  =
  \overline{\bigcup_{\mathcal O\in\mathcal K(M)}\Obs(\mathcal O)}^{\,\|\cdot\|}.
\]
It organizes local observables before a Hilbert-space representation has been
chosen.  Thus the norm closure in \(\Obs_{\rm qloc}\) is part of the abstract
\(C^*\)-algebraic presentation, whereas weak operator closures belong to
represented von Neumann algebras.

\begin{figure}[H]
  \centering
  \resizebox{0.90\textwidth}{!}{%
  \begin{tikzpicture}[x=1cm,y=1cm,every node/.style={font=\scriptsize}]
    \tikzset{
      region/.style={draw,ellipse,line width=0.9pt,fill=blue!5,
        minimum width=2.4cm,minimum height=1.25cm},
      alg/.style={draw,rounded corners=4pt,line width=0.8pt,fill=orange!7,
        minimum width=2.6cm,minimum height=0.75cm,align=center},
      arrow/.style={-{Latex[length=2mm]},line width=0.85pt}
    }
    \node[region,minimum width=3.4cm,minimum height=1.8cm] (o2) at (-3.2,1.0) {\(\mathcal O_2\)};
    \node[region,minimum width=1.7cm,minimum height=0.92cm] (o1) at (-3.2,1.0) {\(\mathcal O_1\)};
    \node[alg] (a1) at (0.65,1.45) {\(\Obs(\mathcal O_1)\)};
    \node[alg] (a2) at (3.85,1.45) {\(\Obs(\mathcal O_2)\)};
    \draw[arrow] (o2.east)--(a1.west);
    \draw[arrow] (a1)--(a2);
    \node at (2.25,1.95) {inclusion};
    \node at (-3.2,-0.05) {\(\mathcal O_1\subset \mathcal O_2\)};
    \node[region] (s1) at (-1.55,-1.25) {\(\mathcal O\)};
    \node[region] (s2) at (2.15,-1.25) {\(\mathcal O'\)};
    \node[alg] (b1) at (-1.55,-2.65) {\(\Obs(\mathcal O)\)};
    \node[alg] (b2) at (2.15,-2.65) {\(\Obs(\mathcal O')\)};
    \draw[arrow] (s1)--(b1);
    \draw[arrow] (s2)--(b2);
    \draw[<->,line width=0.8pt] (b1)--(b2);
    \node[align=center] at (0.3,-3.30) {commuting algebras for spacelike
      separated regions};
  \end{tikzpicture}
  }
  \caption{AQFT assigns algebras to regions; inclusion of regions gives
  inclusion of algebras, and spacelike separation gives commutativity.}
  \label{fig:aqft-local-net}
\end{figure}

\section{Perturbative Algebraic QFT and Interacting Examples}
\label{sec:perturbative-aqft-bridge}

Perturbative algebraic QFT is the framework in which Epstein--Glaser causal
renormalization, local time-ordered products, and algebraic locality are put
in the same language.  It does not by itself construct a nonperturbative
Hilbert-space QFT.  Its objects are formal power series in the interaction
and in \(\hbar\), but the locality and renormalization statements are
mathematically sharp at each order.

For a scalar field on a globally hyperbolic spacetime \(M\), let
\(\mathcal E=C^\infty(M)\) be the classical configuration space.  A
functional \(F:\mathcal E\to\mathbb C\) is called microcausal if each
functional derivative \(F^{(n)}(\varphi)\) is a compactly supported
distribution on \(M^n\) and
\[
  \operatorname{WF}\bigl(F^{(n)}(\varphi)\bigr)
  \cap
  \bigl(\overline V_+^n\cup\overline V_-^n\bigr)
  =
  \emptyset ,
  \label{eq:microcausal-wavefront-condition}
\]
where \(\overline V_\pm^n\) denotes the set of \(n\)-tuples of nonzero
covectors all lying in the closed future or all lying in the closed past
light cone.  This wavefront condition is exactly what makes contractions with
Hadamard two-point functions well-defined.  If \(H(x,y)\) is a chosen
Hadamard bidistribution for the free theory, the free algebra product is
\[
  F\star_H G
  =
  m\circ
  \exp\!\left[
    \hbar
    \int_{M^2}
      H(x,y)
      \frac{\delta}{\delta\varphi(x)}
      \otimes
      \frac{\delta}{\delta\varphi(y)}
      \,\dd x\,\dd y
  \right](F\otimes G),
  \label{eq:paqft-star-product}
\]
where \(m(F\otimes G)(\varphi)=F(\varphi)G(\varphi)\).  The antisymmetric
part of \(H\) is \((\ii/2)\Delta\), with \(\Delta\) the causal propagator, so
the first order of the commutator recovers the Peierls bracket:
\[
  [F,G]_{\star_H}
  =
  \ii\hbar\,\{F,G\}_{\rm Peierls}
  +O(\hbar^2).
\]

The interacting theory is introduced through time-ordered products
\[
  T_n:\mathcal F_{\rm loc}^{\otimes n}\longrightarrow\mathcal F_{\mu{\rm c}},
  \qquad n\ge1,
\]
from local functionals to microcausal functionals.  Away from all diagonals in
configuration space, \(T_n\) is fixed by the free time-ordered contraction.
The Epstein--Glaser problem is the extension across partial diagonals subject
to symmetry, covariance, microlocal spectrum bounds, unitarity, field
independence, and causal factorization.  With
\[
  \mathcal S(F)
  =
  1+\sum_{n\ge1}\frac{\ii^n}{\hbar^n n!}\,T_n(F^{\otimes n}),
\]
causal factorization says
\begin{equation}
  \mathcal S(F+G)=\mathcal S(F)\star_H\mathcal S(G)
  \quad\text{if}\quad
  \operatorname{supp}F\cap J^-(\operatorname{supp}G)=\emptyset .
  \label{eq:paqft-causal-factorization}
\end{equation}
The condition means that the support of \(F\) is not in the causal past of the
support of \(G\), so \(F\) is later than \(G\).  For an interaction functional
\(V\), the relative \(S\)-matrix
\[
  \mathcal S_V(F)
  :=
  \mathcal S(V)^{-1}_{\star_H}\star_H \mathcal S(V+F)
  \label{eq:paqft-relative-s-matrix}
\]
is a formal power series, with the inverse taken in the completed formal
\(\star_H\)-algebra.  The perturbative interacting algebra associated with a
region \(\mathcal O\) is generated by the relative \(S\)-matrices
\(\mathcal S_V(F)\) with \(\operatorname{supp}F\subset\mathcal O\), modulo the
causal factorization relations and the time-slice ideal.  In this sense
perturbative AQFT gives a local net
\[
  \mathcal O\longmapsto \mathfrak A_V^{\rm pert}(\mathcal O)
\]
of formal algebras.

Renormalization in this framework is the freedom in extending \(T_n\) to the
diagonals.  The ambiguity is local: two choices are related by a
Stueckelberg--Petermann redefinition of the interaction by local functionals
and their derivatives.  For gauge theory the same construction must be
combined with the BV complex and the quantum master equation; the physical
perturbative algebra is then the BRST/BV cohomology of the interacting
algebra.  These facts make perturbative AQFT a natural mathematical home for
local counterterms, contact terms, and gauge-theoretic Ward identities.  It
will not be treated only as a dictionary between path integrals and local
nets.  Later chapters must develop the perturbative algebra itself: the
classification of time-ordered-product extensions, the renormalization group
as the Stueckelberg--Petermann group, interacting fields by Bogoliubov's
formula, the time-slice axiom in perturbation theory, BV cohomology for gauge
theories, and the relation between these constructions and the 1PI and
Wilsonian effective actions.

The algebraic formulation must also be developed through substantial
interacting examples.  A persistent weakness of much abstract AQFT literature
is that the axioms, sector theory, and local-algebra technology are developed
with too few nontrivial models in view.  In this monograph, whenever an
interacting construction is available or partially available, the associated
algebraic questions should be asked explicitly: what are the local algebras or
formal local algebras, what is the state space, which sectors are present, how
does scattering or infraparticle structure appear, how are scaling limits and
operator product structures represented, and which steps are theorem-level
rather than expected structure?  Constructive scalar models, perturbatively
constructed Epstein--Glaser/pAQFT models, factorizing two-dimensional models,
lattice-continuum scaling limits, conformal nets, and gauge theories with
BRST/BV control are therefore not examples after the axioms; they are part of
the theory development and tests of whether the algebraic formulation captures
the physics it is meant to organize.

The purpose of including this material is not only to collect known
constructions.  The comparison among Wightman, Euclidean, Wilsonian,
perturbative algebraic, BV, and operator-algebraic formulations exposes
genuine gaps: gauge-theory local nets beyond perturbation theory,
charged sectors requiring nonlocal dressing, infraparticle scattering,
contact-term control in deformations, and the passage from constructive or
lattice scaling limits to detailed algebraic data.  These gaps should be
formulated as precise mathematical and physical problems and, where feasible,
developed within the monograph itself.  Existing results supply boundary
conditions for this development; they are not the endpoint of the project.

\section{States and GNS Representations}

A state on a unital \(C^*\)-algebra \(\Obs\) is a linear functional
\[
  \omega:\Obs\to\mathbb C
\]
such that
\[
  \omega(A^*A)\ge0,\qquad \omega(1)=1.
\]
The Gelfand--Naimark--Segal construction produces a Hilbert-space
representation from \((\Obs,\omega)\).  Let
\[
  \mathcal N_\omega=\{A\in\Obs\mid \omega(A^*A)=0\}.
\]
The positivity of \(\omega\) gives the Cauchy--Schwarz inequality
\[
  |\omega(C^*D)|^2\le \omega(C^*C)\,\omega(D^*D),
  \qquad C,D\in\Obs.
\]
Hence \(C\in\mathcal N_\omega\) implies
\(\omega(C^*D)=\omega(D^*C)=0\) for all \(D\in\Obs\), so that
\(\mathcal N_\omega\) is a linear subspace.  It is also a left ideal.  Indeed,
if \(C\in\mathcal N_\omega\) and \(A\in\Obs\), then the \(C^*\)-algebra order
inequality \(0\le A^*A\le \|A\|^2 1\) gives
\[
  0\le \omega((AC)^*(AC))
  =
  \omega(C^*A^*AC)
  \le
  \|A\|^2\omega(C^*C)
  =
  0 .
\]
Thus \(AC\in\mathcal N_\omega\).  The quotient
\(\Obs/\mathcal N_\omega\) therefore has the well-defined inner product
\[
  \langle [A],[B]\rangle_\omega=\omega(A^*B).
\]
Its completion is \(\Hilb_\omega\).  The representation is
\[
  \pi_\omega(A)[B]=[AB],
\]
which is well-defined because \(\mathcal N_\omega\) is a left ideal.  Moreover
\[
  \|\pi_\omega(A)[B]\|_\omega^2
  =
  \omega(B^*A^*AB)
  \le
  \|A\|^2\omega(B^*B),
\]
so \(\pi_\omega(A)\) extends to a bounded operator on \(\Hilb_\omega\).
The cyclic vector is
\[
  \Omega_\omega=[1].
\]
Then
\[
  \omega(A)=
  \langle\Omega_\omega,\pi_\omega(A)\Omega_\omega\rangle_\omega .
\]

\begin{figure}[H]
  \centering
  \resizebox{0.88\textwidth}{!}{%
  \begin{tikzpicture}[x=1cm,y=1cm,every node/.style={font=\scriptsize}]
    \tikzset{
      box/.style={draw,rounded corners=4pt,line width=0.8pt,fill=blue!4,
        minimum width=3.2cm,minimum height=0.86cm,align=center},
      arrow/.style={-{Latex[length=2mm]},line width=0.85pt}
    }
    \node[box] (alg) at (-4.3,0.85) {abstract algebra\\\(\Obs\)};
    \node[box] (state) at (0,0.85) {state\\\(\omega\)};
    \node[box] (gns) at (4.3,0.85) {GNS triple\\\((\Hilb_\omega,\pi_\omega,\Omega_\omega)\)};
    \node[box,fill=green!7] (local) at (0,-1.25)
      {represented local algebras\\\(\pi_\omega(\Obs(\mathcal O))''\)};
    \draw[arrow] (alg)--(state);
    \draw[arrow] (state)--(gns);
    \draw[arrow] (gns)--(local);
    \draw[arrow] (alg)--(local);
  \end{tikzpicture}
  }
  \caption{A state turns the abstract local net into Hilbert-space operator
  algebras through the GNS construction.}
  \label{fig:gns-representation}
\end{figure}

\section{Vacuum Representations}

Suppose the net is Poincare covariant and \(\omega_0\) is invariant:
\[
  \omega_0(\alpha_g(A))=\omega_0(A).
\]
Then covariance is unitarily implemented in the GNS representation:
\[
  U_{\omega_0}(g)\pi_{\omega_0}(A)U_{\omega_0}(g)^{-1}
  =
  \pi_{\omega_0}(\alpha_g(A)),
  \qquad
  U_{\omega_0}(g)\Omega_{\omega_0}=\Omega_{\omega_0}.
\]
The vacuum spectrum condition is the requirement that the joint spectrum of
the translation generators in this representation lies in
\(\overline V_+\), the closure of the open future light cone \(V_+\) in
momentum space.  A vacuum state is an invariant state satisfying this spectrum
condition and having \(\Omega_{\omega_0}\) as the unique invariant vector up to
phase.

In a vacuum representation, the represented von Neumann algebra of a region is
\[
  \mathcal R_{\omega_0}(\mathcal O)
  =
  \pi_{\omega_0}(\Obs(\mathcal O))''.
\]
The double commutant indicates weak operator closure.  The family
\[
  \mathcal O\mapsto \mathcal R_{\omega_0}(\mathcal O)
\]
is the represented local net; when the vacuum representation is fixed, we often
write \(\mathcal R(\mathcal O)\) for
\(\mathcal R_{\omega_0}(\mathcal O)\).  Thus, in the remainder of this
chapter, \(\Obs(\mathcal O)\) denotes the abstract local \(C^*\)-algebra,
whereas \(\mathcal R(\mathcal O)\) denotes its vacuum-representation von
Neumann completion \(\pi_{\omega_0}(\Obs(\mathcal O))''\); these are related by
the chosen representation and are not interchangeable as notation.

\section{Correlation Functions in a Representation}

The algebraic presentation produces correlation functions once a state and a
representation have been chosen.  If
\[
  A_j\in\Obs(\mathcal O_j),
  \qquad j=1,\ldots,n,
\]
are bounded local observables, their vacuum correlation function is
\[
  \omega_0(A_1\cdots A_n)
  =
  \langle\Omega_{\omega_0},
  \pi_{\omega_0}(A_1)\cdots\pi_{\omega_0}(A_n)
  \Omega_{\omega_0}\rangle .
\]
Covariance gives spacetime dependence by translating the local observables:
\[
  A_j(x_j)=\alpha_{x_j}(A_j).
\]
Thus bounded local observables give well-defined functions or distributions
after smearing in the spacetime parameters \(x_j\), depending on the
continuity and phase-space properties imposed on the net.

Point fields require additional structure.  A field \(\Phi(f)\) localized in
\(\mathcal O\) should be affiliated with the represented algebra
\(\mathcal R(\mathcal O)\), or
should arise as a limit of bounded local observables in a topology specified
by the construction.  When this is available, the Wightman distributions are
recovered as vacuum matrix elements of these affiliated unbounded operators.
The algebraic net therefore supplies the local observable content, while
field coordinatizations specify distributional generators for that content.

\section{Haag Duality and the Dual Net}

For a represented net, the causal complement of a region \(\mathcal O\) is
\[
  \mathcal O^\perp
  =
  \{x\in M:\ (x-y)^2>0\ \text{for all }y\in\mathcal O\}.
\]
Einstein causality gives the inclusion
\[
  \mathcal R(\mathcal O)\subset \mathcal R(\mathcal O^\perp)' .
\]
Haag duality is the equality
\begin{equation}
  \mathcal R(\mathcal O)=\mathcal R(\mathcal O^\perp)' .
  \label{eq:haag-duality}
\end{equation}
It says that every bounded operator commuting with all observables localized
outside \(\mathcal O\) is already generated by observables localized in
\(\mathcal O\).

For any local net, one can form the dual net
\[
  \mathcal R^d(\mathcal O):=\mathcal R(\mathcal O^\perp)' .
\]
When the indexing regions are causally complete,
\(\mathcal O=\mathcal O^{\perp\perp}\), the assignment
\(\mathcal O\mapsto\mathcal R^d(\mathcal O)\) is again local and contains the
original net.  Haag duality is precisely the condition
\(\mathcal R=\mathcal R^d\).  The inclusion
\(\mathcal R\subset\mathcal R^d\) measures the possible gap between the chosen
local observable generators and the maximal algebra determined by causal
commutation.  In superselection theory, Haag duality is the condition that
permits localized representations to be represented cleanly by localized
endomorphisms of the observable algebra.

\section{Point Fields as Coordinatizations}

If a Wightman field \(\Phi(f)\) is available, bounded functions of smeared
fields can generate local algebras.  Conversely, in many algebraic settings
point fields arise as operator-valued distributions affiliated with the local
algebras.  The invariant object is the local observable content:
\[
  \mathcal O\mapsto \Obs(\mathcal O)
  \quad\text{or}\quad
  \mathcal O\mapsto\mathcal R(\mathcal O).
\]
Fields are then coordinatizations of local observables, organized by symmetry,
scaling, and the operator product expansion when those structures are
available.

The passage from unbounded Wightman fields to a Haag--Kastler net is not a
formal consequence of the Wightman axioms in their bare domain formulation.
For real test functions the smeared fields are symmetric operators on a common
invariant domain, but this alone neither supplies self-adjoint closures nor
implies strong commutativity of the spectral measures of spacelike separated
closures.  The construction below records the standard route under the extra
closure and strong-locality hypotheses needed to make it a theorem.

For a self-adjoint operator \(T\) on \(\Hilb\), affiliation with a von
Neumann algebra \(\mathcal M\subset B(\Hilb)\) means that the spectral
projections \(E_T(\Delta)\) of \(T\) belong to \(\mathcal M\) for every Borel
set \(\Delta\subset\mathbb R\).  Equivalently, every unitary in the commutant
\(\mathcal M'\) preserves the graph of \(T\) and commutes with \(T\) on its
domain.  Thus affiliation is the precise way in which an unbounded field can
be localized in a bounded local algebra.

\begin{theorem}[Wightman fields generate a represented local net under strong closure hypotheses]
\label{thm:wightman-to-aqft-net}
Let
\[
  (\Hilb,\vac,U,\mathcal D,\{\widehat\Phi_\alpha\}_{\alpha\in I})
\]
be a Wightman field presentation on \(M\).  Enlarge the field-label space, if
needed, so that the smeared generators used below are Hermitian on
\(\mathcal D\).  Suppose that for every real Hermitian labelled test function
\(h\) the symmetric operator \(\widehat\Phi(h)\) is essentially self-adjoint on
\(\mathcal D\); write \(X(h)\) for its self-adjoint closure and
\[
  E_{h}(\Delta)
\]
for the spectral projection of \(X(h)\) associated with the Borel set
\(\Delta\subset\mathbb R\).  Define, for every region
\(\mathcal O\in\mathcal K(M)\),
\begin{equation}
  \mathcal R_\Phi(\mathcal O)
  :=
  \{E_h(\Delta)\mid
    \operatorname{supp}h\subset\mathcal O,\quad
    h=h^\dagger,\quad
    \Delta\subset\mathbb R\text{ Borel}\}'' .
  \label{eq:wightman-generated-net-from-spectral-projections}
\end{equation}
Equivalently, by the spectral theorem,
\begin{equation}
  \mathcal R_\Phi(\mathcal O)
  =
  \{\exp(\ii t X(h))\mid
    \operatorname{supp}h\subset\mathcal O,\quad
    h=h^\dagger,\quad
    t\in\mathbb R\}'' .
  \label{eq:wightman-generated-net-from-weyl-unitaries}
\end{equation}
Since \(t\) may be absorbed into \(h\), this is the standard displayed
construction
\[
  \mathcal R_\Phi(\mathcal O)
  =
  \{\exp(\ii X(h))\mid
    \operatorname{supp}h\subset\mathcal O,\quad h=h^\dagger\}'' .
\]
For every such \(h\), the closed operator \(X(h)\) is affiliated with
\(\mathcal R_\Phi(\mathcal O)\).  If
\(\mathcal O_1\subset\mathcal O_2\), the same \(X(h)\) is also affiliated with
\(\mathcal R_\Phi(\mathcal O_2)\).
Assume strong locality of these closures: if
\(\mathcal O_1\perp\mathcal O_2\), then every spectral projection entering
\(\mathcal R_\Phi(\mathcal O_1)\) commutes with every spectral projection
entering \(\mathcal R_\Phi(\mathcal O_2)\).  Then
\[
  \mathcal O\longmapsto\mathcal R_\Phi(\mathcal O)
\]
is an isotonic, Poincare-covariant represented local net of von Neumann
algebras on \(\Hilb\).  The vector state
\[
  A\longmapsto\langle\vac,A\vac\rangle
\]
is invariant under the represented Poincare action.
\end{theorem}

\begin{proof}
For a fixed \(h\), the spectral theorem identifies the von Neumann algebra
generated by the projection-valued measure \(E_h\) with the von Neumann
algebra generated by the one-parameter unitary group
\(\{\exp(\ii tX(h))\}_{t\in\mathbb R}\).  Indeed
\(\exp(\ii tX(h))=\int \exp(\ii t\lambda)\,\dd E_h(\lambda)\), so the
unitaries lie in
\(\{E_h(\Delta)\mid\Delta\subset\mathbb R\text{ Borel}\}''\); conversely, the
spectral projections are bounded Borel functions of \(X(h)\), hence lie in
the von Neumann algebra generated by the unitary group by the Borel functional
calculus for a self-adjoint operator.  This proves the equivalence of
\eqref{eq:wightman-generated-net-from-spectral-projections} and
\eqref{eq:wightman-generated-net-from-weyl-unitaries}, and also proves the
affiliation statement because all spectral projections of \(X(h)\) are
generators of \(\mathcal R_\Phi(\mathcal O)\).

If \(\mathcal O_1\subset\mathcal O_2\), then the test functions with support in
\(\mathcal O_1\) are among those with support in \(\mathcal O_2\), hence the
generating spectral projections for
\(\mathcal R_\Phi(\mathcal O_1)\) are contained in
\(\mathcal R_\Phi(\mathcal O_2)\).  This proves isotony.  Wightman covariance
gives
\[
  U(g)X(h)U(g)^{-1}=X(g h),
\]
where \(g h\) is the transformed labelled test function.  The spectral theorem
then gives
\[
  U(g)E_h(\Delta)U(g)^{-1}=E_{g h}(\Delta),
\]
so
\[
  U(g)\mathcal R_\Phi(\mathcal O)U(g)^{-1}
  =
  \mathcal R_\Phi(g\mathcal O).
\]
Strong locality is exactly the commutation of the spectral-projection
generators for spacelike separated regions; taking double commutants preserves
commutation of von Neumann algebras.  The vacuum invariance follows from
\(U(g)\vac=\vac\).
\end{proof}

\begin{remark}[Status of the Wightman-to-net direction]
The theorem is conditional.  Wightman commutativity is initially a statement
about commutators of unbounded smeared fields on a common dense domain; it
does not automatically imply commutativity of the spectral projections of
self-adjoint closures.  Thus the displayed Weyl-unitary formula should not be
read as saying that arbitrary Wightman fields automatically span a local
Haag--Kastler algebra.  In constructive scalar models obtained from Euclidean
measures, the route is model-specific: one proves the Schwinger functions,
OS reconstruction, Wightman properties, and the required local-net properties
with estimates particular to the construction.  That is a proof of the
Haag--Kastler axioms for those models, not a general implication from the
abstract Wightman axioms alone.
\end{remark}

\begin{theorem}[field coordinates on a local net give Wightman data]
\label{thm:aqft-to-wightman-fields}
Let
\[
  \mathcal O\longmapsto\mathcal R(\mathcal O)
\]
be a Poincare-covariant represented local net on a Hilbert space \(\Hilb\),
with unitary implementation \(U\) and invariant unit vector \(\Omega\).  Assume
the translation spectrum of \(U\) is contained in \(\overline V_+\).  Suppose
there are a dense subspace \(\mathcal D\subset\Hilb\), a label set \(I\), and
operators \(\widehat\Phi_\alpha(f)\) such that:
\begin{enumerate}[label=\textup{(\roman*)}]
\item \(\Omega\in\mathcal D\), \(U(g)\mathcal D\subset\mathcal D\), and
\(\widehat\Phi_\alpha(f)\mathcal D\subset\mathcal D\);
\item for each \(\alpha\), the map
\(f\mapsto\widehat\Phi_\alpha(f)\) is an operator-valued tempered
distribution on the common domain \(\mathcal D\);
\item if \(\operatorname{supp}f\subset\mathcal O\), then
\(\widehat\Phi_\alpha(f)\) is affiliated with \(\mathcal R(\mathcal O)\);
\item the fields satisfy the Wightman covariance and adjunction identities on
\(\mathcal D\);
\item fields affiliated with spacelike separated regions have vanishing graded
commutator on \(\mathcal D\);
\item the polynomial field orbit of \(\Omega\) is dense in \(\Hilb\).
\end{enumerate}
Then
\[
  (\Hilb,\Omega,U,\mathcal D,\{\widehat\Phi_\alpha\}_{\alpha\in I})
\]
is a Wightman field presentation, and its vacuum distributions are the matrix
elements
\[
  W_{\alpha_1\cdots\alpha_n}(f_1,\ldots,f_n)
  =
  \langle\Omega,
  \widehat\Phi_{\alpha_1}(f_1)\cdots
  \widehat\Phi_{\alpha_n}(f_n)\Omega\rangle .
\]
\end{theorem}

\begin{proof}
The Hilbert space, invariant vacuum vector, unitary Poincare representation,
dense common domain, and operator-valued tempered distributions are the data
listed in the hypotheses.  The spectral condition is the assumed joint
spectrum condition for translations.  Covariance, adjunction, locality, and
cyclicity are precisely items (iv), (v), and (vi).  Continuity of the vacuum
matrix elements follows from the operator-valued tempered-distribution
hypothesis applied successively on the common invariant domain.  These are the
Wightman field-presentation axioms.
\end{proof}

This relation is important for gauge theories.  The algebraic observable net
can be defined from gauge-invariant observables without choosing gauge
potentials as physical operators.  If charged fields are introduced, they
belong to additional data: an enlarged field algebra, charged
representations, or sector endomorphisms.  They are not elements of the
neutral observable algebra unless the construction explicitly identifies them
as observables.

\section{Reeh--Schlieder and Local Cyclicity}

Let \(\vac\) be the vacuum vector in a Poincare-covariant representation.
Assume the translation spectrum lies in \(\overline V_+\), \(\vac\) is
translation invariant, and the vacuum representation is cyclic for the global
von Neumann algebra
\[
  \mathcal R_{\rm glob}
  :=
  \bigvee_{\mathcal O\in\mathcal K(M)}\mathcal R(\mathcal O).
\]
Here \(\bigvee\) denotes the von Neumann algebra generated by the displayed
family.  For the Reeh--Schlieder theorem, the relevant additivity assumption is
weak additivity, not merely the finite-cover additivity of \(\Obs(\mathcal O)\)
stated above.  In the vacuum representation it is the following
translated-generation condition: for every nonempty open region \(\mathcal U\)
with nontrivial causal completion,
\[
  \bigvee_{a\in\mathbb R^D}\mathcal R(\mathcal U+a)
  =
  \mathcal R_{\rm glob}.
\]
The proof uses exactly this global translated-generation statement.  Ordinary
additivity only identifies the algebra of one region with the algebra generated
by a cover of that same region; it does not by itself say that translates of one
fixed local algebra generate \(\mathcal R_{\rm glob}\).  Under weak additivity,
covariance, global cyclicity, and the spectrum condition, the Reeh--Schlieder
property for any such \(\mathcal O\) is
\[
  \overline{\mathcal R(\mathcal O)\vac}=\Hilb .
\]
Indeed, if \(\psi\) is orthogonal to
\(\mathcal R(\mathcal O)\vac\), choose a nonempty open region
\(\mathcal U\) whose closure is contained in \(\mathcal O\), and choose
\(A\in\mathcal R(\mathcal U)\).  For translations \(a\) in a neighborhood of
\(0\), the inclusion \(\mathcal U+a\subset\mathcal O\) and covariance give
\(U(a)AU(a)^{-1}\in\mathcal R(\mathcal O)\), hence
\[
  \langle \psi,\,U(a)A\vac\rangle=0
\]
there.  The spectrum condition makes this matrix coefficient the boundary value
of a function analytic in the forward tube \(\mathbb R^D+iV_+\), so vanishing
on that open set of real translations forces
\(\langle \psi,U(a)A\vac\rangle=0\) for all real translations \(a\).  Since
\(A\) was arbitrary in \(\mathcal R(\mathcal U)\), weak additivity says that
\(\psi\) is orthogonal to \(\mathcal R_{\rm glob}\vac\).  Global cyclicity then
gives \(\psi=0\).  Thus local operations on the vacuum generate a dense set of
states.  The result is used as an input for modular theory, superselection
analysis, and scattering theory in algebraic form.

\begin{figure}[H]
  \centering
  \resizebox{0.78\textwidth}{!}{%
  \begin{tikzpicture}[x=1cm,y=1cm,every node/.style={font=\scriptsize}]
    \tikzset{
      box/.style={draw,rounded corners=4pt,line width=0.8pt,fill=orange!7,
        minimum width=3.0cm,minimum height=0.86cm,align=center},
      arrow/.style={-{Latex[length=2mm]},line width=0.85pt}
    }
    \node[box] (local) at (-3.6,0.8) {local algebra\\\(\mathcal R(\mathcal O)\)};
    \node[box] (vac) at (0,0.8) {vacuum\\\(\vac\)};
    \node[box] (dense) at (3.6,0.8) {dense subspace\\\(\overline{\mathcal R(\mathcal O)\vac}\)};
    \node[box,fill=green!7] (inputs) at (0,-1.15)
      {covariance + spectrum condition\\+ weak additivity};
    \draw[arrow] (local)--(vac);
    \draw[arrow] (vac)--(dense);
    \draw[arrow] (inputs)--(dense);
  \end{tikzpicture}
  }
  \caption{The Reeh--Schlieder property expresses local cyclicity of the
  vacuum in positive-energy local QFT.}
  \label{fig:reeh-schlieder-cyclicity}
\end{figure}

\section{Locally Covariant Form}

On a general globally hyperbolic spacetime \(X\), locality can be expressed as
a functorial assignment
\[
  X\longmapsto \Obs(X).
\]
An isometric causal embedding \(\chi:X\to Y\) induces a homomorphism
\[
  \alpha_\chi:\Obs(X)\to\Obs(Y)
\]
compatible with composition.  The time-slice property says that if \(\chi(X)\)
contains a Cauchy surface of \(Y\), then \(\alpha_\chi\) is an isomorphism.
This formulation preserves the local content of QFT without relying on a
distinguished global vacuum in spacetimes where no such state is selected by
symmetry.

For the present monograph, the algebraic framework supplies a precise language
for local observables, states, sectors, and structural locality theorems.  Its
connection with path integrals, perturbative expansions, and point-field
descriptions is a comparison problem to be tracked explicitly case by case.
